\chapter{Remaining Work}
The current prototype implements repository analysis, system-model construction, capability realization, bounded parameter retrieval, structured language-model inference, deterministic validation, and template-based generation. It also includes a 150-mission paraphrase dataset and two prompt experiments. The remaining work therefore concerns evaluation quality and deployment evidence rather than initial implementation.

The planned activities are as follows:

\begin{itemize}
    \item Complete independent author review of all generated paraphrases and freeze development and held-out partitions before further prompt changes.
    \item Add independently annotated unsupported and ambiguous missions so that unsafe acceptance, unsupported-request rejection, and clarification recall can be measured.
    \item Separate retrieval, parameter-selection, and value-reasoning errors through context ablations and an oracle-selection experiment that supplies the correct identifiers to the value stage.
    \item Repeat the principal experiments to estimate run-to-run variation and report confidence intervals, since temperature zero alone does not guarantee identical inference.
    \item Evaluate the second semantic-enrichment schema against the deterministic context variants, particularly whether aliases, user expressions, and typed relationships improve retrieval and coupled-parameter selection.
    \item Refine prompts or compare alternative local models only on the development partition, then evaluate the frozen candidate once on the untouched held-out set.
    \item Exercise complete generated bundles in ROS~2 simulation and, where feasible, on the physical AntRobot platform, including behavioral checks, failure handling, and reproducibility across repeated deployments.
\end{itemize}

These activities will determine whether the structurally complete prototype also provides reliable semantic behavior. In particular, they address the main limitation exposed by the current results: partial capability understanding does not compensate for poorly grounded parameter selection in an end-to-end configuration task.


% \section{Indicații formatare tabele}
% Se recomandă utilizarea tabelelor de forma celui de mai jos.  Font size :  9. 
% Orice tabel prezent în teză va fi referit în text; exemplu: a se vedea Tabel~\ref{tab:criterii}.

% \begin{table}[th]\small\linespread{1}
% \caption{Sumarizare criterii}
% \label{tab:criterii}
% \begin{tabular}{l >{\raggedright\arraybackslash}p{8cm} >{\raggedright\arraybackslash}p{4cm}}
% \textbf{Calificativ} & \textbf{Criteriu} & \textbf{Observații} \\\hline
% \textbf{Nesatisfacator} & Sunt prezentate pe scurt scheme și pseudo-cod & \\\hline
% \textbf{Satisfacator} &Descriere sumara a implementării, prezentarea unor secvențe nerelevante de cod, scheme, etc.& \\
% \hline
% \textbf{\textit{Bine}} &Descrierea detaliată a algoritmilor/structurilor utilizați; Prezentarea etapizată a dezvoltării, inclusiv cu dificultăți de implementare întâmpinate, soluții descoperite; (dacă este cazul) demonstrarea corectitudinii algoritmilor utilizați. & Pot fi incluse configurații, secvente de cod, pseudo-cod, implementări ale unor algoritmi, analize ale unor date, scripturi de testare. \\
% \hline
% \end{tabular}
% \end{table}
